\documentclass{article}
\usepackage[T1]{fontenc}
\usepackage{lmodern}
\usepackage{graphicx}
\usepackage{amsmath,amssymb}
\usepackage[a4paper,margin=2cm]{geometry}
\usepackage[absolute,overlay]{textpos}
\setlength{\TPHorizModule}{1mm}
\setlength{\TPVertModule}{1mm}
\usepackage[indonesian]{babel}
\usepackage{hyperref}
\setlength{\emergencystretch}{2em}
% unit: Halaman 60

\begin{document}

% === Halaman 60 (python/native) ===

60

Kode Program RC4 (dalam Bahasa Python) 

from Crypto.Cipher import ARC4

import base64

def RC4\_encrypt(plainteks : str, kunci : bytes) -> str:

    cipher = ARC4.new(kunci) 

    cipherteks = cipher.encrypt(plainteks.encode())

    return base64.b64encode(cipherteks).decode()

def RC4\_decrypt(cipherteks : str, kunci : bytes) -> str:

    cipher = ARC4.new(kunci) 

    cipherteks\_decoded = base64.b64decode(cipherteks)

    return cipher.decrypt(cipherteks\_decoded)

pesan = input("Ketikkan pesan rahasia:")

kunci = input("Ketikkan kata kunci: ")

kunci = bytearray(kunci, 'utf-8')

ciphertext = RC4\_encrypt(pesan, kunci)

print(f"Cipherteks: \{ciphertext\}")

plaintext = RC4\_decrypt(ciphertext, kunci)

print(f"Plainteks semula: \{plaintext\}")

\end{document}
